The demographic shift from the Rust Belt to the Sun Belt between 2020 and 2030 is
the largest cross-decade reapportionment since 1990. Texas alone projects to gain
three seats (38 to 41); New York projects to lose two. These changes require new
bisection tree designs and reveal algorithmic challenges specific to fast-growth
suburban geography.

Three findings define the Sun Belt redistricting challenge.
First, TX $k = 41$ is prime: the bisect-apportion binary fallback produces a
20+21 top-level split, which is algorithmically sound and empirically validated
in this paper.
Second, Sun Belt growth is concentrated in suburban rings (Houston exurbs, Phoenix
Valley, Atlanta suburbs), creating new design challenges where the metro-suburb
bisection is the natural first cut but the suburb is growing faster than the city.
Third, the alternative GeoSection algorithm (\texttt{--structure ratio-optimal})
handles any $k$ via ratio-optimal bisection without requiring prime factorization,
and is recommended as the preferred 2030 structure for TX and FL given their
irregular prime factorizations.

For declining states, the main algorithmic concern is large-scale district mergers:
NY losing two seats requires merging or substantially redesigning one-seventh of
the 2020 congressional map. Algorithmic plans provide a systematic, auditable merger
process that enacted gerrymanders cannot.
